% !TEX TS-program = xelatex
% Generated from exams/g11_midterm_v1.json. No external TeX inputs required.
\documentclass[11pt,a4paper]{article}
\usepackage{fontspec,xeCJK}
% Use the engine default Latin Modern Roman for portability.
\IfFontExistsTF{Noto Serif CJK TC}{\setCJKmainfont{Noto Serif CJK TC}}{
  \IfFontExistsTF{Songti TC}{\setCJKmainfont{Songti TC}}{\setCJKmainfont{FandolSong-Regular.otf}}}
\usepackage[margin=17mm,top=18mm,bottom=20mm]{geometry}
\usepackage{enumitem,tabularx,array,fancyhdr,lastpage,needspace}
\setlength{\parindent}{0pt}
\setlength{\parskip}{2pt}
\setlength{\emergencystretch}{2em}
\linespread{1.08}
\setlength{\headheight}{14pt}
\pagestyle{fancy}
\fancyhf{}
\fancyhead[L]{\small G11-U13-01}
\fancyhead[R]{\small 學生試卷}
\fancyfoot[C]{\small 第\thepage 頁／共\pageref{LastPage}頁}
\renewcommand{\headrulewidth}{0pt}
\renewcommand{\footrulewidth}{0pt}
\newenvironment{qblock}{\par\vspace{3pt}\noindent\begin{minipage}{\linewidth}}{\end{minipage}\par\vspace{4pt}}
\newcommand{\smallheading}[1]{\par\vspace{6pt}{\large\bfseries #1}\par\smallskip}
\begin{document}
\begin{center}{\LARGE 高二英文練習段考}\par
{\large Unit 13 與綜合能力練習}\end{center}
班級：\underline{\hspace{1.5cm}}　座號：\underline{\hspace{1cm}}　姓名：\underline{\hspace{2.2cm}}\par
範圍：Core Vocabulary Unit 13、基礎文法與篇章閱讀、資訊整合、中譯英及短文寫作
\par\smallskip
\begin{enumerate}[leftmargin=1.5em,nosep]\small
\item 全卷共六頁、三十六題，滿分一百分，作答時間七十分鐘。請直接在試卷上作答，不使用答案卡。
\item 第一、二、四大題及第二十九至三十題為四選一，每題僅選一個答案，填在題號前括號內。第三大題填字庫的字母。
\item 第三十一至三十六題依指示以英文作答。選擇題多選、未作答或修改後無法辨認者不給分；答錯不倒扣。
\item 請勿使用字典、翻譯工具或其他電子設備。閱讀材料為命題用原創情境。
\end{enumerate}
\begin{center}\small\begin{tabular}{|c|c|c|c|c|c|c|c|}\hline
詞彙 & 綜合 & 選填 & 閱讀 & 混合 & 翻譯 & 寫作 & 總分 \\ \hline
２０ & １０ & １０ & １６ & １４ & １０ & ２０ & １００ \\ \hline
\rule{0pt}{18pt} & & & & & & & \\ \hline\end{tabular}\end{center}\section*{一　詞彙與搭配}
{\small 第１至１０題，每題２分。依句意選出最適當的單字。
\par\smallskip
}
\begin{qblock}
\textbf{(\hspace{0.6cm}) 1.} The club cannot afford new cameras because its annual \_\_\_\_\_ is only three thousand dollars.
\par
\begin{tabularx}{\linewidth}{@{}XXXX@{}}
(A) report & (B) survey & (C) budget & (D) policy
\end{tabularx}
\end{qblock}
\begin{qblock}
\textbf{(\hspace{0.6cm}) 2.} The lawyer's \_\_\_\_\_ paid her to represent him in court and explain the legal process.
\par
\begin{tabularx}{\linewidth}{@{}XXXX@{}}
(A) client & (B) colleague & (C) judge & (D) assistant
\end{tabularx}
\end{qblock}
\begin{qblock}
\textbf{(\hspace{0.6cm}) 3.} Excited about her first trip abroad, Amy was \_\_\_\_\_ to learn a few useful phrases before leaving.
\par
\begin{tabularx}{\linewidth}{@{}XXXX@{}}
(A) afraid & (B) unable & (C) ashamed & (D) eager
\end{tabularx}
\end{qblock}
\begin{qblock}
\textbf{(\hspace{0.6cm}) 4.} To remove the extra water, \_\_\_\_\_ the wet sponge firmly over the sink.
\par
\begin{tabularx}{\linewidth}{@{}XXXX@{}}
(A) stretch & (B) squeeze & (C) scratch & (D) slide
\end{tabularx}
\end{qblock}
\begin{qblock}
\textbf{(\hspace{0.6cm}) 5.} All members must \_\_\_\_\_ their votes before the meeting ends; late votes will not be counted.
\par
\begin{tabularx}{\linewidth}{@{}XXXX@{}}
(A) cast & (B) draw & (C) hold & (D) earn
\end{tabularx}
\end{qblock}
\begin{qblock}
\textbf{(\hspace{0.6cm}) 6.} Please \_\_\_\_\_ the muddy footprints off the floor with a damp cloth before anyone slips.
\par
\begin{tabularx}{\linewidth}{@{}XXXX@{}}
(A) pour & (B) fold & (C) wipe & (D) hang
\end{tabularx}
\end{qblock}
\begin{qblock}
\textbf{(\hspace{0.6cm}) 7.} When the dry curtains caught fire, tall \_\_\_\_\_ quickly rose toward the ceiling.
\par
\begin{tabularx}{\linewidth}{@{}XXXX@{}}
(A) clouds & (B) flames & (C) shadows & (D) sparks
\end{tabularx}
\end{qblock}
\begin{qblock}
\textbf{(\hspace{0.6cm}) 8.} Still \_\_\_\_\_ about being unfairly blamed, Eric refused to speak to the teammate who had accused him.
\par
\begin{tabularx}{\linewidth}{@{}XXXX@{}}
(A) curious & (B) nervous & (C) cheerful & (D) bitter
\end{tabularx}
\end{qblock}
\begin{qblock}
\textbf{(\hspace{0.6cm}) 9.} The child would answer questions only in the \_\_\_\_\_ of her mother, who had to remain beside her.
\par
\begin{tabularx}{\linewidth}{@{}XXXX@{}}
(A) absence & (B) distance & (C) presence & (D) silence
\end{tabularx}
\end{qblock}
\begin{qblock}
\textbf{(\hspace{0.6cm}) 10.} Hold both ends of the wet towel and \_\_\_\_\_ them in opposite directions to force the water out.
\par
\begin{tabularx}{\linewidth}{@{}XXXX@{}}
(A) twist & (B) shake & (C) raise & (D) press
\end{tabularx}
\end{qblock}

\newpage
\section*{二　綜合測驗}
{\small 第１１至１５題，每題２分。閱讀短文，選出各空格的最佳答案。
\par\smallskip
}
\textbf{A Repair Afternoon}\par\smallskip
Last month, our school held a repair afternoon to give broken everyday objects a second chance. Students \_\_\_11\_\_\_ had never fixed anything before worked with adult volunteers at tables set up inside the library. \_\_\_12\_\_\_ some volunteers expected students to lose interest quickly, the young helpers stayed until the last visitor left and asked to join the next event. By the time the doors opened at two, the adults \_\_\_13\_\_\_ three lamps to demonstrate simple repairs and had placed them near the entrance. Each student had \_\_\_14\_\_\_ part in a short safety lesson before being allowed to use any tools. The experience increased their \_\_\_15\_\_\_: several students said they now believed they could solve small problems at home without immediately buying replacements. Objects that could not be fixed that day were labeled clearly, so their owners knew what kind of professional help to seek. The organizers hoped that the event would become a regular part of school life.
\par\smallskip
\begin{qblock}
\textbf{(\hspace{0.6cm}) 11.} 
\par
\begin{tabularx}{\linewidth}{@{}XXXX@{}}
(A) which & (B) who & (C) whose & (D) where
\end{tabularx}
\end{qblock}
\begin{qblock}
\textbf{(\hspace{0.6cm}) 12.} 
\par
\begin{tabularx}{\linewidth}{@{}XXXX@{}}
(A) Because & (B) Unless & (C) After & (D) Although
\end{tabularx}
\end{qblock}
\begin{qblock}
\textbf{(\hspace{0.6cm}) 13.} 
\par
\begin{tabularx}{\linewidth}{@{}XXXX@{}}
(A) had repaired & (B) will repair & (C) have repaired & (D) are repairing
\end{tabularx}
\end{qblock}
\begin{qblock}
\textbf{(\hspace{0.6cm}) 14.} 
\par
\begin{tabularx}{\linewidth}{@{}XXXX@{}}
(A) made & (B) kept & (C) taken & (D) given
\end{tabularx}
\end{qblock}
\begin{qblock}
\textbf{(\hspace{0.6cm}) 15.} 
\par
\begin{tabularx}{\linewidth}{@{}XXXX@{}}
(A) curiosity & (B) confidence & (C) patience & (D) sympathy
\end{tabularx}
\end{qblock}
\section*{三　文意選填}
{\small 第１６至２０題，每題２分。從字庫選出最適當的詞，填入字母；每詞最多用一次，有三個多餘選項。
\par\smallskip
}
\begin{center}\begin{tabular}{llll}
(A) adventure & (B) burden & (C) passion & (D) clue\\
(E) routine & (F) reward & (G) rush & (H) scream\\
\end{tabular}\end{center}
\textbf{A Mystery Walk}\par\smallskip
Our drama club created a mystery walk through the old part of our town. Mei's \_\_\_16\_\_\_ for local history was clear: she spent every free afternoon reading old maps simply because she loved the subject. The team wanted visitors to feel a sense of \_\_\_17\_\_\_ as they explored unfamiliar lanes and discovered places they had never noticed before. When a cardboard ghost suddenly fell from a window, one visitor let out a loud \_\_\_18\_\_\_ that everyone in the next street could hear, then laughed in relief. A sudden \_\_\_19\_\_\_ of visitors arrived just after lunch, so the team introduced starting times to prevent crowding. At the final stop, a muddy footprint became the decisive \_\_\_20\_\_\_ that helped visitors identify which character had entered the garden. After finishing the walk, families compared their ideas over cups of tea and explained which details had changed their minds. The club kept their comments to improve the next event.
\par\smallskip
Answers: 16. \_\_\_\_\_\_   17. \_\_\_\_\_\_   18. \_\_\_\_\_\_   19. \_\_\_\_\_\_   20. \_\_\_\_\_\_
\par\smallskip

\newpage
\section*{四　閱讀測驗}
{\small 第２１至２８題，每題２分。依各篇文章選出最佳答案。
\par\smallskip
}
\textbf{Reading A　The Cup Return Problem}\par\smallskip
At Brook High School, students could borrow a reusable cup from the cafe by paying a small deposit. They received the money back when they returned the cup. The student council expected this system to reduce waste, but many cups remained in classrooms at the end of each day. In interviews, students said they usually remembered the cups; the problem was that the cafe closed before their last class ended.
\par\smallskip
For the next two weeks, the council placed return boxes near the school gates and kept the deposit unchanged. The number of cups returned on the day they were borrowed rose from 60 out of 100 to 85 out of 100. The cafe also washed every returned cup before lending it again. However, the council did not count how many disposable cups students bought during either period.
\par\smallskip
Some members wanted to announce that the project had reduced waste by 25 percent. Others disagreed. A higher return rate showed that returning borrowed cups had become easier, but it did not reveal how much total waste the school produced. The council therefore decided to record both cup returns and purchases of disposable cups during a longer trial. It also planned to ask students who still kept their cups overnight what prevented them from returning them.
\par\smallskip
{\small deposit＝押金；disposable＝拋棄式的。
\par\smallskip
}
\begin{qblock}
\textbf{(\hspace{0.6cm}) 21.} What was the main reason for adding return boxes near the gates?
\par
\begin{enumerate}[label=(\Alph*),leftmargin=1.0cm,nosep,topsep=2pt]
\item To make students pay a larger deposit.
\item To let students wash cups after class.
\item To stop the cafe from lending more cups.
\item To let students return cups after the cafe closed.
\end{enumerate}
\end{qblock}
\begin{qblock}
\textbf{(\hspace{0.6cm}) 22.} Which statement is directly supported by the recorded numbers?
\par
\begin{enumerate}[label=(\Alph*),leftmargin=1.0cm,nosep,topsep=2pt]
\item More borrowed cups were returned on the same day.
\item Students bought fewer disposable cups during the trial.
\item The school's total waste fell by exactly one quarter.
\item Every student returned a cup before going home.
\end{enumerate}
\end{qblock}
\begin{qblock}
\textbf{(\hspace{0.6cm}) 23.} In paragraph 2, what does `either period' refer to?
\par
\begin{enumerate}[label=(\Alph*),leftmargin=1.0cm,nosep,topsep=2pt]
\item The first and last classes on the same school day.
\item The cafe's opening hours and the time after it closed.
\item The periods before and after the boxes were added.
\item The time before and after a borrowed cup was washed.
\end{enumerate}
\end{qblock}
\begin{qblock}
\textbf{(\hspace{0.6cm}) 24.} Why did the council plan a longer trial with another kind of record?
\par
\begin{enumerate}[label=(\Alph*),leftmargin=1.0cm,nosep,topsep=2pt]
\item To prove that every unreturned cup had been lost.
\item To judge the environmental result with more complete evidence.
\item To replace student interviews with guesses about their habits.
\item To find a reason to remove all the new return boxes.
\end{enumerate}
\end{qblock}

\newpage
\section*{四　閱讀測驗（續）}
\textbf{Reading B　From Writing to Speaking}\par\smallskip
In the first month of an English discussion club, twelve students attended regularly, but only a few spoke. The leader, Lina, first thought the quiet members were not interested. Instead of requiring everyone to speak immediately, she asked each member to write one reason for staying silent. Several students said they needed time to organize their ideas. Others worried that a mistake would interrupt the discussion.
\par\smallskip
Lina changed the next four meetings. Before each discussion, members spent two minutes writing. They then shared their ideas with one partner before speaking to the whole group. Lina also allowed members to submit a question on paper. By the fourth meeting, nine students had contributed at least once, compared with four at the first meeting. One student, however, said that speaking with the same partner every week limited the range of opinions he heard.
\par\smallskip
The club decided to keep the preparation time but change partners regularly. Lina did not claim that the new activities alone explained the improvement: members might also have become more familiar with one another. Still, their written comments helped her understand a problem she had first misread. For her, participation was no longer simply a matter of asking quiet people to try harder. It also depended on how the activity was organized.
\par\smallskip
\begin{qblock}
\textbf{(\hspace{0.6cm}) 25.} Which title best expresses the main idea?
\par
\begin{enumerate}[label=(\Alph*),leftmargin=1.0cm,nosep,topsep=2pt]
\item Changing a Discussion So More Members Can Join In
\item Winning a Speaking Contest with Perfect Grammar
\item Choosing the Most Experienced Leader for a Club
\item Replacing Group Discussions with Written Tests
\end{enumerate}
\end{qblock}
\begin{qblock}
\textbf{(\hspace{0.6cm}) 26.} Which change directly addressed the need to organize ideas before speaking?
\par
\begin{enumerate}[label=(\Alph*),leftmargin=1.0cm,nosep,topsep=2pt]
\item Requiring every member to speak as soon as a meeting began.
\item Allowing only the most active members to ask questions.
\item Keeping the same discussion partner for the whole year.
\item Giving members two minutes to write before sharing.
\end{enumerate}
\end{qblock}
\begin{qblock}
\textbf{(\hspace{0.6cm}) 27.} Why did the club decide to change partners regularly?
\par
\begin{enumerate}[label=(\Alph*),leftmargin=1.0cm,nosep,topsep=2pt]
\item Quiet students had asked to stop all pair work.
\item A member wanted to hear a wider variety of views.
\item Lina wanted to avoid reading written comments.
\item The club no longer had time for preparation.
\end{enumerate}
\end{qblock}
\begin{qblock}
\textbf{(\hspace{0.6cm}) 28.} Why was Lina careful about explaining the increase in participation?
\par
\begin{enumerate}[label=(\Alph*),leftmargin=1.0cm,nosep,topsep=2pt]
\item No one had written down how many students spoke.
\item All members had already mastered public speaking.
\item Growing familiarity might also have helped students speak.
\item The changes had made every member less willing to talk.
\end{enumerate}
\end{qblock}

\newpage
\begingroup\small\section*{五　混合題}
{\small 第２９至３２題，每題２分；第３３題６分。閱讀公告與表格，再依各題指示作答。
\par\smallskip
}
\textbf{School Service Day}\par\smallskip
Registration closes at 4:00 p.m. on Friday. Each student may join ONE team. Check the schedule and requirements before choosing. The school will lend all necessary tools. Students must bring a water bottle, but no previous volunteer experience is required. Work in pairs and follow the team leader's instructions throughout the activity.
\par\smallskip
Outdoor teams will move to the gym if it rains; their starting and ending times will stay the same. The Office Team will work in the school office in all weather. Students must attend the full session, including the final cleanup. A bus to Central Station leaves the school at 11:40 a.m. Students who need it should choose a team that finishes before then.
\par\smallskip
{\small closed-toe shoes＝包覆腳趾的鞋。
\par\smallskip
}
\begin{center}\small\renewcommand{\arraystretch}{1.15}\begin{tabularx}{\linewidth}{|l|l|X|l|}\hline
Team & Time & Main task & Special requirement\\ \hline
Garden Team & 09:00-11:00 & Plant flowers outdoors & Bring a hat\\ \hline
Office Team & 10:00-12:00 & Sort donated books indoors & Read small labels\\ \hline
Repair Team & 09:30-11:30 & Fix chairs outdoors & Wear closed-toe shoes\\ \hline
\end{tabularx}\end{center}
Leo cannot arrive before 9:20 a.m. He needs to take the 11:40 a.m. bus, can read small labels, and owns closed-toe shoes. He has no previous volunteer experience.
\par\smallskip
\begin{qblock}
\textbf{(\hspace{0.6cm}) 29.} Which team can Leo join while meeting all his time limits?
\par
\begin{tabularx}{\linewidth}{@{}XXXX@{}}
(A) Garden Team only & (B) Office Team only & (C) Garden Team or Repair Team & (D) Repair Team only
\end{tabularx}
\end{qblock}
\begin{qblock}
\textbf{(\hspace{0.6cm}) 30.} Which statement correctly describes the arrangements for a rainy day?
\par
\begin{enumerate}[label=(\Alph*),leftmargin=1.0cm,nosep,topsep=2pt]
\item The Repair Team changes location but keeps its original hours.
\item The Garden Team starts later and no longer needs a full session.
\item The Office Team moves to the gym with all the outdoor teams.
\item All teams finish early so every student can catch the bus.
\end{enumerate}
\end{qblock}
\begin{qblock}
\textbf{31.} Copy TWO consecutive words from the notice to complete the sentence: Every student must bring a \_\_\_\_ \_\_\_\_.
\par\smallskip
\end{qblock}
\begin{qblock}
\textbf{32.} Copy TWO consecutive words from the notice: On a rainy day, outdoor teams keep the same starting and \_\_\_\_ \_\_\_\_.
\par\smallskip
\end{qblock}
\begin{qblock}
\textbf{33.} Sam writes, `I have never volunteered before, so I cannot join any team.' In ONE or TWO English sentences, correct his misunderstanding and tell him ONE action he must complete before Service Day. Use the registration deadline in your answer.
\par\smallskip
\par\smallskip\noindent\vbox{
\vskip 0.65cm\hrule width\linewidth height0.2pt
\vskip 0.65cm\hrule width\linewidth height0.2pt
\vskip 0.65cm\hrule width\linewidth height0.2pt
}\par
\end{qblock}
\endgroup
\newpage
\section*{六　中譯英}
{\small 第３４至３５題，每題５分。將完整句意譯成自然、正確的英文，不限定使用單一句型。
\par\smallskip
}
\begin{qblock}
\textbf{34.} 這場雨使我們無法在戶外舉行原定於今天下午進行的校園音樂會。
\par\smallskip
\par\smallskip\noindent\vbox{
\vskip 0.65cm\hrule width\linewidth height0.2pt
\vskip 0.65cm\hrule width\linewidth height0.2pt
}\par
\end{qblock}
\begin{qblock}
\textbf{35.} 你越仔細閱讀這些說明，完成這項任務時就會犯越少錯誤。
\par\smallskip
\par\smallskip\noindent\vbox{
\vskip 0.65cm\hrule width\linewidth height0.2pt
\vskip 0.65cm\hrule width\linewidth height0.2pt
}\par
\end{qblock}
\section*{七　短文寫作}
{\small 第３６題，２０分。依提示寫一封８０至１００字的英文電子郵件。
\par\smallskip
}
內容８分、組織４分、語言６分、拼寫與標點２分。問候語及署名不計字數。少於７０或超過１１０字另扣１分。
\par\smallskip
36. Your class has received NT\$3,000 to improve English learning. Write an email to your class representative, Alex, suggesting ONE activity. Explain what students will do, how the money will be used, why the activity will help, and ONE possible problem with a practical solution. Do not sign your real name; use Chris.
\par\smallskip
\par\smallskip\noindent\vbox{
\vskip 0.65cm\hrule width\linewidth height0.2pt
\vskip 0.65cm\hrule width\linewidth height0.2pt
\vskip 0.65cm\hrule width\linewidth height0.2pt
\vskip 0.65cm\hrule width\linewidth height0.2pt
\vskip 0.65cm\hrule width\linewidth height0.2pt
\vskip 0.65cm\hrule width\linewidth height0.2pt
\vskip 0.65cm\hrule width\linewidth height0.2pt
\vskip 0.65cm\hrule width\linewidth height0.2pt
\vskip 0.65cm\hrule width\linewidth height0.2pt
\vskip 0.65cm\hrule width\linewidth height0.2pt
\vskip 0.65cm\hrule width\linewidth height0.2pt
\vskip 0.65cm\hrule width\linewidth height0.2pt
\vskip 0.65cm\hrule width\linewidth height0.2pt
\vskip 0.65cm\hrule width\linewidth height0.2pt
}\par
Word count: \_\_\_\_\_\_
\par\smallskip

\end{document}
